\documentclass[12pt]{article}
\usepackage[utf8]{inputenc}
\usepackage[T1]{fontenc}
\usepackage{hyperref}
\title{Turumsa Language Resources for K-12 Curriculum Development}
\author{Compiled by Claude Sonnet 5 (claude-sonnet-5) for Andr\'as Kornai}
\date{}
\begin{document}
\maketitle

\section{Introduction}
Turumsa (ISO 639-3 tqm) is a possibly extinct language of Makapa village, Gogodala Rural LLG, Middle Fly District, Papua New Guinea. Only five elderly speakers were found in 2002. Most Makapa residents now speak Dibiyaso. Its closest relative is Doso, at about 61\% lexical similarity, and the two are grouped as Doso--Turumsa. The classification is uncertain. The survey brief calls it Trans-New Guinea, and other sources treat it as unclassified or as a Bosavi language.

\section{Primary Educational Resources}
None exist as far as the searches showed.

\section{Linguistic and Reference Materials}
\begin{itemize}
\item Wikipedia, Turumsa language: \url{https://en.wikipedia.org/wiki/Turumsa_language}
\item Wikipedia, Doso--Turumsa languages: \url{https://en.wikipedia.org/wiki/Doso%E2%80%93Turumsa_languages}
\item Wikipedia, Doso language: \url{https://en.wikipedia.org/wiki/Doso_language}
\end{itemize}

\section{Religious and Community Resources}
None were found.

\section{Digital and Audio Resources}
The only other entry found is a catalogue record, with no audio or text content seen.
\begin{itemize}
\item Endangered Languages Project, Turumsa: \url{https://endangeredlanguages.com/lang/3512}
\end{itemize}

\section{Curriculum Development Context}
With so few remaining speakers and no known documentation beyond brief descriptions, any curriculum work would be a revitalisation or documentation effort. It would have to begin with recording the remaining speakers. Further searching is needed for any word lists or field data.
\end{document}